%======================================================================
\section{The zero-killing family and Theorem D}\label{sec:family}
%======================================================================

Every function of the form $F=\Xi^2H$, with $H$ analytic on a strip containing $\lvert\Im t\rvert\le\frac12$, vanishes at every
non-trivial zero of $\zeta$, on or off the critical line. We call such functions \emph{zero-killing}. For such $F$ the explicit formula loses its zero side, and $\Arch(F)$
becomes a sum over the prime powers (Corollary~\ref{cor:zerokilling}). Section~\ref{ssec:bessel} gives the Fourier transform of $\Xi^2H$
in closed form; Section~\ref{ssec:hermite} defines the functions $\Xi^2P_J(t^2)$ whose Fourier transforms have zeros of order at least two at the
first $J$ prime powers; and Section~\ref{ssec:thmM} proves Theorem~\ref{thm:intro-D}.

Here $\xi(s)=\frac12s(s-1)\GR(s)\zeta(s)$ with $\GR(s)=\pi^{-s/2}\Gamma(s/2)$, $\Xi(t)=\xi(\frac12+it)$ is Riemann's $\Xi$-function, and
$\Xiinf(t):=\frac12s(s-1)\GR(s)$ at $s=\frac12+it$ is its archimedean part. We use the coordinates $x=e^{2\pi\xi}$ and $y=2\pi x$, so that
the prime power $n$ sits at $\xi=\xi_n$, at $x=n$ and at $y=2\pi n$, and $\theta=x\,\frac{d}{dx}=y\,\frac{d}{dy}=\frac1{2\pi}\frac{d}{d\xi}$.
For $a<\pi/2$, $\Hc_a$ is the class of even functions $H$, real on $\RR$ and analytic on some closed strip
$\lvert\Im t\rvert\le\frac12+\delta$ ($\delta>0$), with $\lvert H(t)\rvert\le Ce^{a\lvert\Re t\rvert}$ there; $\Hc:=\bigcup_{a<\pi/2}\Hc_a$.

%----------------------------------------------------------------------
\subsection{The Bessel form}\label{ssec:bessel}
%----------------------------------------------------------------------

We use two classical facts about $\Xi$. First, Riemann's representation
$\Xi(t)=\int_\RR\Phi(u)e^{iut}\dd u$ with
$\Phi(u)=2\sum_{n\ge1}\bigl(2\pi^2n^4e^{9u/2}-3\pi n^2e^{5u/2}\bigr)\exp(-\pi n^2e^{2u})$, where
$\Phi$ is even \cite[p.~305, (2)--(3) and the remark after (4)]{Pol26}. Second, Stirling's formula and the convexity bound for
$\zeta$ give, for every $b\in(0,1]$, a constant $C_b$ with
\begin{equation}\label{eq:XB}
\lvert\Xi(t)\rvert\le C_b(1+\lvert\Re t\rvert)^{3}e^{-\pi\lvert\Re t\rvert/4}\qquad(\lvert\Im t\rvert\le b)
\end{equation}
\cite{Tit86}. (The same bound holds for $b\le\frac52$; only $b<1$ is used.) By~\eqref{eq:XB}, $\Xi^2H\in\Tc_\delta$ whenever $H\in\Hc_a$ is analytic near
$\lvert\Im t\rvert\le\frac12+\delta$ (the decay $e^{-(\pi/2-a)\lvert\Re t\rvert}$ absorbs any power of $1+\lvert z\rvert$);
in particular $\Xi^2P(t^2)\in\Tc_\delta$ for every polynomial $P$ and every $\delta\in(0,\frac12)$.

\begin{lemma}[Positivity of the kernel]\label{lem:Psi-pos}
$\Psi(\xi)=\widehat{\Xi^2}(\xi)>0$ for every $\xi\in\RR$.
\end{lemma}

\begin{proof}
For $u\ge0$ every term of the series for $\Phi$ is positive, because $2\pi n^2e^{2u}\ge2\pi>3$; as
$\Phi$ is even, $\Phi>0$ on $\RR$ (see \cite[Theorem~A(i)]{CNV86}, which credits P\'olya \cite{Pol27}). Hence
$\widehat\Xi(\xi)=2\pi\Phi(2\pi\xi)>0$, and $\Psi=\widehat\Xi*\widehat\Xi>0$.
\end{proof}

Write $K_0,K_1$ for the modified Bessel functions and
\begin{equation}\label{eq:W0}
\Wnull(z):=z^2\bigl[(z^2+9)K_0(z)-6zK_1(z)\bigr]=\theta^2(\theta+1)^2K_0(z),\qquad \theta=z\tfrac{d}{dz}.
\end{equation}
The second expression follows from $\theta K_0=-zK_1$ and $\theta(zK_1)=-z^2K_0$: on the basis
$e_0=K_0$, $e_1=zK_1$ one has $(\theta+1)^2e_0=(1+z^2)e_0-2e_1$, then
$\theta(\theta+1)^2e_0=4z^2e_0-(1+z^2)e_1$, and finally
$\theta^2(\theta+1)^2e_0=z^2(z^2+9)e_0-6z^2e_1$.

\begin{proposition}[Bessel form of $\Psi$]\label{thm:bessel}
For every real $\xi$, with $x=e^{2\pi\xi}$,
\[
\Psi(\xi)=\int_\RR\Xi(t)^2e^{-2\pi i\xi t}\dd t=2\pi\sqrt x\sum_{N\ge1}d(N)\,\Wnull(2\pi Nx),
\]
the series converging absolutely and uniformly on compact sets, together with all its
$\xi$-derivatives. In particular $\Psi$ is even and real, and
$\Psi(\xi_n)=(C_\infty+o(1))\,n^4e^{-2\pi n}$ as $n\to\infty$, with
$C_\infty=\pi(2\pi)^4=(2\pi)^{9/2}(\pi/2)^{1/2}=4896.3\ldots$.
The Gamma-only analogue
$\Psione(\xi):=\widehat{\Xiinf^{\,2}}(\xi)=2\pi\sqrt x\,\Wnull(2\pi x)$ also holds, and extends to complex
$x$ with $\lvert\arg x\rvert<\pi/2$.
\end{proposition}

Here $\Xiinf$ is not real on $\RR$, so $\Psione$ is real but not even, and $\Xiinf^{\,2}P(t^2)$ is not a test function; $\Psione$ enters only
through Lemma~\ref{lem:abel}. The proof is in Appendix~\ref{app:proofs-bessel}. One has
$\int_0^\infty\sum_Nd(N)K_0(2\pi N\sqrt\lambda)\,\lambda^{s/2-1}\dd\lambda=\frac12\GR(s)^2\zeta(s)^2$; in the variable $u$ with $\lambda=e^{2u}$, the operator $(\partial_u^2-\frac14)^2$ removes the double poles
at $s=0$ and $s=1$, and Fourier inversion identifies the result with $\Xi^2$.

\begin{remark}[Attribution]
Proposition~\ref{thm:bessel} is a differentiated form of the extended Koshliakov formula. Koshliakov's
formula for $\sum_Nd(N)K_0(2\pi Nx)$ \cite{Kos29} had been found earlier by Ramanujan
\cite[Ch.~3]{AB13}; its extended form, with an integral of $\Xi^2$, is due to Dixit \cite{Dix10}, and in
\cite[Theorem~1.4]{Dix12} it expresses $\sqrt x\sum_Nd(N)K_0(2\pi Nx)$, up to an explicit polar term, as a
cosine transform of $\Xi(t/2)^2/(1+t^2)^2$. The operator $(\theta^2-\frac14)^2$ removes both the polar
term and the factor $(t^2+\frac14)^{-2}$. For more general Koshliakov-kernel identities of this type see
\cite[Theorem~1.2]{DRRZ16}. We include the proof (Appendix~\ref{app:proofs-bessel}) because the
normalisation matters for the certificates of Section~\ref{sec:certified}.
\end{remark}

\begin{corollary}[The explicit formula for zero-killing functions]\label{cor:zerokilling}
Let $F=\Xi^2H$ with $H$ analytic on $\lvert\Im t\rvert\le\frac12+\delta$, and assume $F\in\Tc$ (for
instance $H\in\Hc$, or $H$ a polynomial in $t^2$). Then
\[
\Arch(F)=\frac1\pi\sum_{n\in\PP}\Lambda(n)n^{-1/2}\FT F(\xi_n).
\]
In particular, if $\FT F$ vanishes at $\xi_{n_1},\dots,\xi_{n_J}$, then
$\Arch(F)=\frac1\pi\sum_{n\in\PP,\,n>n_J}\Lambda(n)n^{-1/2}\FT F(\xi_n)$. No hypothesis on the zeros
of $\zeta$ is used.
\end{corollary}

\begin{proof}
Apply Lemma~\ref{lem:explicit-formula}. For every non-trivial zero $\rho$, on or off the line,
$t_\rho=-i(\rho-\frac12)$ satisfies $\lvert\Im t_\rho\rvert<\frac12$ and
$\Xi(t_\rho)=\frac12\rho(\rho-1)\GR(\rho)\zeta(\rho)=0$, so $F(t_\rho)=\Xi(t_\rho)^2H(t_\rho)=0$ and
the zero side vanishes.
\end{proof}

%----------------------------------------------------------------------
\subsection{The exact family}\label{ssec:hermite}
%----------------------------------------------------------------------

Let $n_1<n_2<\cdots$ be the prime powers, $\PP=\{n_j\}$, and put $m_k:=\widehat{t^{2k}\Xi^2}$; since $(P(t^2)F)^\wedge=P(-\theta^2)\FT F$,
$m_k=(-\theta^2)^k\Psi$. All derivatives below are taken in $\xi$.

\begin{proposition}[The exact family is a linear Hermite interpolation]\label{prop:E}
Let $J\ge1$. If the matrix $M_J=[m_k(\xin{n_j});m_k'(\xin{n_j})]_{j\le J,\,1\le k\le2J}$ is non-singular, that is $J\in\Jset$, there is exactly one
polynomial $P_J$ with $P_J(0)=1$ and $\deg P_J\le2J$ such that $\FT F_J:=\widehat{\Xi^2P_J(t^2)}$ has zeros of
order at least two at $\xin{n_1},\dots,\xin{n_J}$. Its coefficient vector is $p^{(J)}=-M_J^{-1}b_J$, with
$b_J=[m_0(\xin{n_j});m_0'(\xin{n_j})]$. Equivalently, $\FT F_J$ is the error of Hermite interpolation of
$\Psi$ by $\operatorname{span}\{m_1,\dots,m_{2J}\}$ at the nodes. Each $F_J=\Xi^2H_J$, $H_J(t)=P_J(t^2)$,
lies in $\Tc$, and $\Arch(F_J)=\frac1\pi\sum_{n\in\PP,\,n>n_J}\Lambda(n)n^{-1/2}\FT F_J(\xi_n)$.
\end{proposition}

\begin{proof}
We have $\FT F_P=\sum_kp_km_k$ for $F_P=\Xi^2P(t^2)$, so the $2J$ conditions are the
linear system $M_J(p_1,\dots,p_{2J})^{\mathsf T}=-b_J$. The element $-\sum_{k\ge1}p_k^{(J)}m_k$ of
$\operatorname{span}\{m_1,\dots,m_{2J}\}$ has the data of $m_0=\Psi$, which is the second description. The rest is
the remark before Lemma~\ref{lem:Psi-pos} and Corollary~\ref{cor:zerokilling}.
\end{proof}

Every real polynomial with $P(0)=1$ factors over $\RR$ as $\prod_j(1+r_ju+s_ju^2)$, by pairing complex
conjugate roots and pairing the real roots arbitrarily (with one factor having $s_j=0$ if the degree is odd). The cushioned certificates of an earlier version
of this paper, used as a cross-check in Section~\ref{ssec:kappa} (ancillary directory \nolinkurl{kappa/}), use this parametrisation of $P_J$.

\begin{lemma}[An integral representation]\label{lem:abel}
For a polynomial $P$ put $\mathcal Q_P(\theta):=P\bigl(-(\theta+\tfrac12)^2\bigr)\theta^2(\theta+1)^2$ and
$Q_P(y):=e^{y}\,\mathcal Q_P(\theta)\,e^{-y}$, with $\theta=y\frac{d}{dy}$. Then $Q_P$ is a polynomial in $y$ of degree $2\deg P+4$, with
leading coefficient $(-1)^{\deg P}$ times that of $P$, and $Q_P(0)=0$. For $z>0$,
\[
g_P(z):=\bigl[P\bigl(-(\theta+\tfrac12)^2\bigr)\Wnull\bigr](z)=\int_z^\infty Q_P(y)e^{-y}(y^2-z^2)^{-1/2}\dd y,
\]
and $\widehat{\Xi^2P(t^2)}(\xi)=2\pi\sqrt x\sum_Nd(N)g_P(2\pi Nx)$.
\end{lemma}

The proof is in Appendix~\ref{app:proofs-hermite}. The certificates of Proposition~\ref{prop:exactcone} (Appendix~\ref{app:exact}) evaluate $\widehat{\Xi^2P_J(t^2)}$ and its derivatives through Proposition~\ref{thm:bessel} and this lemma.

%----------------------------------------------------------------------
\subsection{Robust compactness and the proof of Theorem D}\label{ssec:thmM}
%----------------------------------------------------------------------

The next lemma is the compactness step behind
every conditional theorem below. It needs no convergence of the family and no exact positivity at finite $J$.

\begin{lemma}[Robust compactness]\label{lem:robust}
Let $H_J$, $J$ in an infinite index set, be even, real on $\RR$ and analytic on $\lvert\Im t\rvert\le\frac12+\delta$,
and put $F_J:=\Xi^2H_J$. Assume:
\begin{enumerate}[label=(\roman*),leftmargin=2em]
\item $\lvert H_J(t)\rvert\le Ce^{a\lvert\Re t\rvert}$ on $\lvert\Im t\rvert\le\frac12+\delta$, with $C$ and $a<\pi/2$ independent of $J$;
\item $H_J(0)=1$;
\item $\liminf_JH_J(t)\ge0$ for every real $t$;
\item $\liminf_J\FT F_J(\xi)\ge0$ for every $\xi\ge\xi_2$;
\item $\FT F_J(\xi_n)\to0$ for every prime power $n$.
\end{enumerate}
Then a subsequence $H_J$ converges locally uniformly on $\{\lvert\Im t\rvert<\frac12+\delta\}$ to some $H_\infty$, and
$F_\infty:=\Xi^2H_\infty\in\Cone\cap\Tc$, with $\Arch(F_\infty)=0$ and $\FT F_\infty(0)=\int F_\infty>0$. Hence
$\kstar\le0$; moreover $\Arch(F)\ge0$ for every $F\in\Cone$ of the form $\Xi^2H$, so the infimum of
$\Arch(F)/\int F$ over such $F$ is $0$ and is attained at $F_\infty$.
\end{lemma}

The proof (Appendix~\ref{app:proofs-M}) combines Montel's theorem with dominated convergence: by (i) and
\eqref{eq:XB}, $\lvert F_J(t)\rvert\le C(1+\lvert\Re t\rvert)^6e^{-(\pi/2-a)\lvert\Re t\rvert}$ on the strip, uniformly in $J$.

Conditions (iv) and (v) are
asymptotic, so cushioned or rationalised representatives may be used, provided the uniform bound (i)
holds for them. The hypotheses of Lemma~\ref{lem:robust}, for some sequence, are equivalent to the existence of an
exact magic function of the form $\Xi^2H$ with $H\in\Hc$ and $H(0)\ne0$ (take the constant sequence $H/H(0)$), which is a priori
stronger than \condS.

For the exact family $(P_J)_{J\in\Jset}$ of Proposition~\ref{prop:E} we consider the following hypotheses.
\begin{itemize}[leftmargin=3.2em]
\item[\hyp{N}] $\Jset$ is infinite.
\item[\hyp{G$_a$}] There are $\delta\in(0,\frac12)$, $a<\pi/2$ and $C$ with $\lvert H_J(t)\rvert\le Ce^{a\lvert\Re t\rvert}$ on
$\lvert\Im t\rvert\le\frac12+\delta$ for all $J\in\Jset$.
\item[\hyp{Z$_\infty$}] No real point is a limit of zeros of $H_{J_k}$ with $J_k\in\Jset$, $J_k\to\infty$.
\item[\hyp{Mg}] For every $\xi\ge\xi_2$, $\liminf_{J\in\Jset}\FT F_J(\xi)\ge0$.
\end{itemize}
Given \hyp{G$_a$}, \hyp{Z$_\infty$} is equivalent to: for every $T>0$,
$\liminf_{J\in\Jset}\min_{\lvert t\rvert\le T}H_J(t)>0$ (Appendix~\ref{app:proofs-M}). It holds, for instance, if there are
$\eta>0$ and $J_0$ such that every zero of $H_J$, $J\in\Jset$, $J\ge J_0$, has $\lvert\Im t\rvert\ge\eta$. Since a bound on a strip implies
the same bound on every narrower strip, the size of $\delta$ in \hyp{G$_a$} plays no role. Finite-$J$ instances of \hyp{N}, and of the
coefficient bound \eqref{eq:POSa} below, which would give \hyp{G$_a$} and \hyp{Z$_\infty$}, are certified for $J\in\{10,60,61,110,111\}$
(Proposition~\ref{thm:finiteJ}); \hyp{Mg} is a $\liminf$ in $J$ and has no finite-$J$ instance, although its finite-$J$ analogue
$\FT F_J\ge0$ on $[\xin2,\infty)$ is certified at the same $J$ (Proposition~\ref{prop:exactcone}).

\begin{theorem}[Theorem D]\label{thm:criterion}
Assume \hyp{N}, \hyp{G$_a$}, \hyp{Z$_\infty$} and \hyp{Mg} for the exact family. Then some subsequential limit
$F_\infty=\Xi^2H_\infty$, with $H_\infty\in\Hc$ and $H_\infty>0$ on $\RR$, satisfies $F_\infty\in\Cone\cap\Tc$,
$\Arch(F_\infty)=0$ and $F_\infty^{-1}(0)\cap\RR=\Zz$. Hence \condS\ and \condU\ hold, and $\RH\Leftrightarrow\condE$; under either,
$\kstar=0$ and $\Kad=\{\pz\}$.
\end{theorem}

In an earlier version of this paper \hyp{Mg} asked for $\liminf_{\Jset}\FT F_J(\xi)>0$ at every $\xi\in[\xi_2,\infty)\setminus\xiPP$, so that
the zeros of $\FT F_\infty$ in $[\xi_2,\infty)$ lay in $\BRS$. Theorem~\ref{thm:zerosupport} makes that unnecessary; this weakening is due to
Astra (OpenAI), contributed during an independent review of an earlier version of this paper, and independently verified. Under the
strict form the proof below gives, in addition, $\FT F_\infty^{-1}(0)\cap[\xi_2,\infty)=\xiPP$.

\begin{proof}
Hypotheses (i), (ii) of Lemma~\ref{lem:robust} are \hyp{G$_a$} and $P_J(0)=1$. For (v): if $n=n_j$, then
$\FT F_J(\xi_n)=0$ for every $J\in\Jset$ with $J\ge j$. Hypothesis (iv) is \hyp{Mg}. For (iii): by the equivalent form of
\hyp{Z$_\infty$}, for each $T$ there is $c_T>0$ with $H_J\ge c_T$ on $[-T,T]$ for all large $J\in\Jset$. So Lemma~\ref{lem:robust} applies,
and its limit satisfies $H_\infty\ge c_T>0$ on $[-T,T]$ for every $T$. Hence $F_\infty^{-1}(0)\cap\RR=\Zz$. (Under the strict form of
\hyp{Mg}, $\FT F_\infty(\xi)=\lim\FT F_J(\xi)>0$ for $\xi\in[\xi_2,\infty)\setminus\xiPP$, which gives the additional statement above.)
Lemma~\ref{lem:robust} gives \condS, and Corollary~\ref{cor:magic}(a), applied to $F_\infty$, gives \condU, and RH under \condE; conversely RH
implies \condE{} (Theorem~\ref{thm:logic}). Under either, $\kstar\ge0$ (Theorem~\ref{thm:duality}), so
$\kstar=0$, and $\Kad=\{\pz\}$.
\end{proof}

No convergence of the family is assumed, and the proof never divides by $\Xi^2$: $H_\infty$ inherits the bound
\hyp{G$_a$}, so $F_\infty\in\Tc$ directly. The strict inequality $a<\pi/2$ is essential for the domination. Further
details (norm convergence, the equivalent forms of \hyp{Z$_\infty$}) are in Appendix~\ref{app:proofs-M}. In sphere packing, Cohn and Miller
conjectured, with numerical evidence, that explicit sequences of functions converge to the magic functions in dimensions $8$ and $24$
\cite{CM16}; Theorem~\ref{thm:criterion} asks for the same kind of convergence, for an explicit family, in the explicit formula.

Hypotheses \hyp{G$_a$} and \hyp{Z$_\infty$} concern the polynomials $P_J$ alone. They follow from a bound on the coefficients.

\begin{proposition}[Coefficient bounds control the zero side]\label{prop:coeff}
Suppose that for all $J\in\Jset$ with $J\ge J_0$,
\begin{equation}\label{eq:POSa}
0\le p_k^{(J)}\le C\,\frac{a^{2k}}{(2k)!}\quad\text{for all }k,\quad\text{with }0\le a<\pi/2\text{ and }C\text{ independent of }J.
\end{equation}
Then for these $J$: (1) $H_J(t)\ge1$ for all real $t$, so \hyp{Z$_\infty$} holds; (2) $\lvert H_J(t)\rvert\le C\cosh(a\lvert t\rvert)$ for all $t\in\CC$, so \hyp{G$_a$} holds on every strip;
(3) $\{H_J\}$ is normal on $\CC$, and every limit $H_\infty$ is entire of exponential type at most $a$, with
$0\le[t^{2k}]H_\infty\le Ca^{2k}/(2k)!$ and $H_\infty\ge1$ on $\RR$; (4) $\lvert P_J(u)\rvert\le P_J(\lvert u\rvert)$, and every root $u$ of
$P_J$ satisfies $\lvert u\rvert\ge\min_kp_{k-1}^{(J)}/p_k^{(J)}$ (minimum over $k$ with $p_k^{(J)}>0$).
\end{proposition}

The proof is in Appendix~\ref{app:proofs-zeroside}.

\begin{remark}[Variants]\label{rem:variant}
(a) In Theorem~\ref{thm:criterion}, \hyp{Z$_\infty$} may be replaced by a lower bound $H_J\ge h$ on $\RR$ for all large $J\in\Jset$, with
$h:\RR\to(0,\infty)$; the pointwise form $\liminf_{\Jset}H_J(t)>0$ for every real $t$ suffices, and the limit then satisfies $H_\infty\ge h>0$.
By Proposition~\ref{prop:coeff}(1), the coefficient bound \eqref{eq:POSa} for all large $J$ gives this with $h\equiv1$; so \hyp{N}, \hyp{Mg}
and \eqref{eq:POSa} for all large $J$ imply \condS\ and \condU.
(b) If $\kstar=0$ and some weak magic function $F^*$ has only finitely many real zeros outside $\Zz$ and satisfies $\widehat Z(F^*)\subseteq\BRS$, then
RH holds and $\Kad=\{\pz\}$: $\kstar=0$ gives \condE\ (Theorem~\ref{thm:duality}), by Theorem~\ref{thm:crit}(c) every admissible pair is
carried by $Z(F^*)\times\widehat Z(F^*)$, and Theorem~\ref{thm:Gfin} applies to each of them. If $Z(F^*)\subseteq\Zz\cup\{0\}$, the condition
on $\widehat Z(F^*)$ is not needed: every admissible pair is then carried by $\Zz\cup\{0\}$ on the zero side, and
Theorem~\ref{thm:zerosupport} applies.
\end{remark}

A second route to \condU\ and \condS\ uses certificates directly, without a limit function.

\begin{corollary}[A certificate route]\label{cor:certroute}
Suppose that there are $F_J\in\Cone$, for $J$ in an infinite set, and constants $c_J>0$ with $F_J\ge c_J\,\Xi^2$ on $\RR$ and
$\Arch(F_J)/c_J\to0$. Then \condU\ and \condS\ hold.\astranote
\end{corollary}

\begin{proof}
For $(\mu,\nu)\in\Kad$, Lemma~\ref{lem:weak-duality} gives $0\le c_J\int\Xi^2\dd\mu\le\int F_J\dd\mu\le\Arch(F_J)$, so $\int\Xi^2\dd\mu=0$; as
$\Xi^2>0$ off $\Zz$, $\mu$ is carried by $\Zz$, and Theorem~\ref{thm:zerosupport} gives $(\mu,\nu)=\pz$. Moreover $\int F_J\ge c_J\int\Xi^2>0$, so
$\kstar\le\Arch(F_J)/\int F_J\le\max(0,\Arch(F_J)/c_J)/\int\Xi^2\to0$.
\end{proof}

The route needs exact cone members $F_J$ at infinitely many $J$; when $\Arch(F_J)>0$ for every $J$, as for the certificates of
Section~\ref{ssec:kappa}, no finite ladder of certificates can supply it. It needs no compactness and no \hyp{Z$_\infty$}. We do not know
whether either set of hypotheses implies the other. The certificates of Section~\ref{ssec:kappa} have
$c_J=1$, and give the finite-$J$ bound of Corollary~\ref{cor:nearrigid}.
